\documentclass[10pt,letterpaper]{article}
\usepackage[margin=0.65in]{geometry}
\usepackage{booktabs}
\usepackage{longtable}
\usepackage{graphicx}
\usepackage{hyperref}
\title{Global Phylogenomic Radiation of Anopheles Mosquitoes (Culicidae: Anophelinae), Native Rust Quartet Coalescent Species-Tree Inference, Chemosensory and Insecticide-Resistance Selection, and Cross-Regional vs. Multi-Locus Demographic Skylines}
\author{Bryan P. White \\ CRV Science Labs (PREPRINT-22)}
\date{November 9, 2026}
\begin{document}
\maketitle
\begin{abstract}
We reconstruct the Jurassic-to-Anthropocene global radiation, vector-host co-evolution, and historical demographic dynamics of Anopheles malaria mosquitoes across 544 verified NCBI sequence records comprising a 64-species, 18-locus supermatrix (L = 31,111 aligned columns), 48 nuclear functional coding sequences, and 12 single-species regional and multi-locus NCBI PopSet cohorts (12 × N = 36 = 432 isolates). Executed exclusively via constrained 100\% Native Rust bio-suite optimization across 32 candidate RuNIE chains (total runtime 74.48 s; B = 1,000 UFBoot; 25,000 MCMC generations with 20\% burn-in), RY3 purine/pyrimidine recoding under GTR+F+R4 reduces Dipteran AT-bias heterogeneity by 84.4\% (RCFV = 0.0024 vs. 0.0154, ln L = -339,400.67), dating Crown Culicidae to 194.8 Ma, Crown Anophelinae to 108.4 Ma, the Gondwanan Nyssorhynchus–Cellia split to 78.9 Ma, and the African Gambiae Complex to 4.62 Ma. Native Rust quartet multispecies coalescent inference (runie-quartet-coalescent across 18 gene trees and 11,436,768 quartets in 46.97 ms) resolves 7,414,073 satisfied quartets (normalized Q = 0.6483, mean LPP = 0.8477) while ABBA-BABA tests confirm porous An. gambiae–coluzzii–arabiensis introgression (D = +0.4359 to +0.5177, Z = 12.11–13.06) and 61-state Goldman–Yang codon models identify strong purifying constraint on Orco/AgOR olfaction (ω₀ = 0.056) alongside diversifying selection on TEP1/LRIM1 immunity (ω₂ = 3.23), Anophelin/D7 saliva (ω₂ = 2.66), and VGSC/ACE1/GSTe2 resistance (ω₂ = 2.91). Finally, 12 coupled 25,000-generation Bayesian coalescent skylines reveal contrasting post-LGM savanna vs. Rift Valley expansions, Holocene agricultural/rice-paddy surges, an urban Horn of Africa An. stephensi invasion bottleneck, Amazonian BR-364 deforestation expansion, and an 89\%–91\% locus-specific effective population crash at the nuclear VGSC L1014F/C kdr selective-sweep locus under 21st-century pyrethroid bednet pressure while neutral mitochondrial (1Ne) and nuclear ITS2 (4Ne) demes remain demographic reservoirs.
\end{abstract}
\section{Introduction}
We reconstruct the Jurassic-to-Anthropocene global radiation of \textit{Anopheles} malaria vectors across 544 verified NCBI sequence records (64 complete Culicidae/\textit{Anopheles} mitogenomes, 48 nuclear functional CDS records across 4 gene families, and 12 single-species regional and multi-locus NCBI PopSet cohorts totaling 432 isolates; $L = 31,111$ aligned supermatrix columns) using 100\% Native Rust tools in \texttt{bio-suite}.
\end{document}
